\section{The schedule}
\label{sec:10.2}
Table 10.1 says, for a removal pipeline, what Table 17.1 says for targeting: it says what the instruments yield, and what was considered and rejected. It is also the schedule of the statute in \S\ref{sec:10.4}. Chapter~\ref{sec:4} took the instruments entire so that any omission would show. The entries below come from running that test against what the 2025–26 operations did, as Amnesty documented them \autocite{Amnesty2026Whistles} and as \S\ref{sec:10.1} records.

\runin{Table 10.1. The removal schedule}

\textbf{UDHR 14; Refugee Convention Art. 33 (via the Protocol); CAT Art. 3, read with General Comment 4 \autocite{cat2017gc4}.} Don't take part in removing anyone whose claim to protection hasn't been finally decided. Don't take part in removing anyone to a place where they face persecution or torture, directly or through a third country that may send them on. Don't take part in transferring anyone into a foreign government's custody where they would be held without charge or contact. A receiving government's assurance doesn't satisfy this entry. Rejected: don't process protection claims at all, because a claim has to be assessed to be granted. Rejected: accept diplomatic assurances, because the Committee against Torture treats them as no substitute for assessing the person's own risk \autocite{cat2017gc4}.

\textbf{UDHR 14; Refugee Convention Art. 33; CAT Art. 3; ICCPR 13.} Don't take part in removing anyone who hasn't been asked, in a language they understand, whether they fear return, and told how to claim protection. Rejected: rely on the person to raise a fear, because the people least likely to know the words are the ones expedited removal moves fastest (\S\ref{sec:10.3}, position 6).

\textbf{UDHR 13(2); ICCPR 12(4), read with General Comment 27 \autocite{hrc1999gc27}.} Don't take part in removing anyone from a country that is their own, including long-term residents with deep ties there. Rejected: citizens only, because the Human Rights Committee reads “own country” more broadly. The reading is General Comment 27's and is contested in the Committee's own cases, which have gone both ways on long-term residents \autocite{hrc1996stewart,hrc2011nystrom}; the entry is a departure and takes effect as one. (Chapter~\ref{sec:4} sets the bar a departure must clear.)

\textbf{UDHR 12, 16(3); ICCPR 17, 23, 24.} Don't take part in a removal that separates a parent from a minor child. Rejected: don't hold family records, because they are needed to keep families together. Rejected: avoid separation by detaining the family together, because the next entry bars it.

\textbf{ICCPR 24, read with General Comment 35 \autocite{hrc2014gc35} and the joint general comment on children in migration \autocite{cmwcrc2017jgc}.} Don't take part in detaining a child for immigration purposes, alone or with family. Rejected: detention of children only as a last resort and for the shortest appropriate time, because that is General Comment 35's reading, and what an agency says of every detention it orders. The entry follows the joint comment, which treats detaining a child because of the child's or the parents' migration status as always contrary to the child's best interests; it is a departure and takes effect as one. (Chapter~\ref{sec:4} sets the bar a departure must clear.)

\textbf{UDHR 9, 10; ICCPR 9, 13.} Don't take part in a detention or removal made without an individual determination by a person, on the person's own case. A determination made at a rate no adjudicator could sustain didn't occur, and neither did one that consists of a score, a points table or a membership designation. Rejected: a model-side term requiring review, because the model can't establish whether review occurred.

\textbf{UDHR 3, 5; ICCPR 6, 7, read with General Comment 36 \autocite{hrc2018gc36}.} Don't take part in force beyond what is strictly necessary to make an arrest, or in lethal force except where strictly unavoidable to protect life. Don't mark a person as dangerous, armed or gang-affiliated on a designation, field or score the person hasn't had a chance to contest. Rejected: don't record danger at all, because officers are owed warning of a real threat. The entry bars uncontested designations, not information.

\textbf{UDHR 5; ICCPR 7, 10; CAT Arts. 1, 16.} Don't take part in holding anyone in conditions that are cruel, inhuman or degrading, or incommunicado. Don't take part in a placement or transfer that puts a detained person out of reach of counsel or family. Rejected: leave conditions to inspection, because an inspection reported to the operator is the operator's record.

\textbf{Refugee Convention Art. 31.} Don't take part in penalizing a refugee for unlawful entry where they came directly and presented themselves.

\textbf{UDHR 2, 7; ICCPR 2, 26.} Don't use race, ethnicity, religion or national origin, or proxies for them such as language, accent, apparent ethnicity, location or type of work, as inputs to selecting whom to stop, detain or remove. Rejected: don't record them, because auditing for discriminatory effect needs them.

\textbf{UDHR 19; ICCPR 19.} Don't use a person's opinions or speech, published or private, as a ground for detention, removal or revoking status. Don't screen speech to find such grounds. Rejected: screen speech only for support of designated groups, because the screening of 2025 was found to have selected people for their political speech \autocite{aaup2025}.

\textbf{UDHR 20; ICCPR 19, 21, 22, read with General Comment 37 \autocite{hrc2020gc37}.} Don't take part in identifying, tracking or selecting anyone for observing, recording or assisting people during enforcement, or for association with those who do. Rejected: protect accredited monitors only, because the people who kept the record in 2025–26 were neighbours, and nobody accredited them.

\textbf{UDHR 12, 20; ICCPR 17, 21, read with General Comment 37.} Don't take part in identifying people by their face or other biometrics in public places for enforcement. Rejected: bar biometrics in custody too, because identification at booking, with notice, is how a person's own records are found and corrected.

\textbf{UDHR 12; ICCPR 17, 26.} Don't use data collected to provide health care, collect taxes, educate or deliver social services to identify, locate or select anyone for enforcement. Rejected: allow use under a court order, because the January 2026 Medicaid transfer went beyond what a court had allowed \autocite{JoffeBlock2026MedicaidPalantir}.

These terms cover most removals now under way. They don't cover every removal. A person who has been asked whether they fear return, and who has no protection claim and no tie to the country, can still be removed after an individual determination. That most of what the agency now does fails these terms is the case for ending it. The own-country entry and the children's entry are departures and take effect only as departures.

Several entries answer something specific in the record. The force entry reaches a model because the pipeline's outputs set how agents approach a door. A dossier field that codes someone as gang-affiliated is securitization in one line, the person in need described as the person to fear. Amnesty documented killings by federal agents during these operations, some of which it says may amount to extrajudicial executions, and detention conditions it says should be investigated as torture \autocite{Amnesty2026Whistles}.

The race entry names proxies because of what happened in Los Angeles. A district court had barred immigration stops there based on apparent race or ethnicity, speaking Spanish or accented English, location and type of work, and in September 2025 the Supreme Court stayed that order \autocite{vasquezperdomo2025}. An entry that named race alone would allow the combination the order had barred.

The observers' entry protects the record this chapter relies on. Amnesty's report takes its title from them, \emph{We Had Whistles, They Had Guns}. It found that video taken by rapid responders and community members was critical to establishing what happened in several of the most serious incidents \autocite{Amnesty2026Whistles}. A pipeline that can be pointed at the people keeping the record will be, and the record goes next. Amnesty asks the government to end unlawful surveillance and the use of facial recognition for identification, and to protect the rights to peaceful assembly, expression and association \autocite{Amnesty2026Whistles}. The observers' and biometric entries are those asks in the schedule's form.

The last entry, the firewall, rests on privacy and equality, which is the ground the floor's instruments supply. Its strongest justification lies outside them. People who fear that a clinic visit, a tax return or a school enrolment will be used to find them stay away, and what they lose is health care, education and income support. Those are economic and social rights, which Chapter~\ref{sec:4} left out of the floor. The entry belongs on the floor because it protects the floor's rights. It belongs in the statute because it also protects the others, and a statute can carry what the floor can't (\S\ref{sec:10.4.5}).
